\section{Evaluation}\label{sec:tsps-evaluation}

To assess the concrete performance of our protocols, we implement and
benchmark the schemes of Section~\ref{sec:tsps-abe}
and~\ref{sec:tsps-kiltz}. We also describe two applications of efficient
threshold structure-preserving signatures: transparency logs, and digital
assets on a permissioned ledger.

\paragraph{Transparency logs.} Certificate transparency (CT) and key
transparency (KT) are the main examples of transparency logs.\footnote{A
citation for CT and KT is still to be added.} In both cases, a transparency
log is an append-only, publicly verifiable record, of either certificates or
keys. Each entry is sent to a log provider, signed, and appended to the log;
in CT systems, the log server returns a Signed Certificate Timestamp (SCT) to
the certificate provider, and in KT systems, the mapping from user identity to
public key is what the log server signs. Transparency logs let users and
auditors verify that the certificates or keys they interact with carry
legitimacy beyond the app serving them, moving trust onto the log providers
instead. Threshold signatures could remove this single point of trust,
distributing it across several servers: a threshold structure-preserving
signature given to the key or certificate owner would let the
structure-preserving property of the signature be used to prove, in
zero-knowledge, that a key is recorded in the log, for instance for
privacy-preserving identity verification. Since the structure-preserving
property lets the public key itself be signed directly, without first
hashing it, the resulting zero-knowledge proofs are more efficient than they
would otherwise be.

\paragraph{Permissioned ledgers.} A permissioned ledger is a distributed
ledger where token issuance, transaction verification, and participation are
controlled by one or a small number of authorities. Permissioned ledgers are
of particular interest for platforms such as Central Bank Digital Currencies
(CBDCs), as they offer a transparent system enabling regulatory oversight
while still supporting transaction privacy and compliance with the monetary
policy of the central bank in question.\footnote{A citation for CBDC
permissioned ledger designs is still to be added.} One way to build such a
system has every transaction executed and, if valid, signed by an endorser
before being added to the ledger,\footnote{A citation for endorser-based
ledger designs is still to be added.} where the signed state committed to the
ledger is a commitment to the value and type of asset. Threshold
structure-preserving signatures would let trust be distributed between
multiple endorsers via the threshold property, while the structure-preserving
property lets users perform efficient proofs of statements about their
assets on the ledger, for example for compliance purposes.

\paragraph{Realising the ideal MPC functionalities.} Section~\ref{sec:MPC}
describes the ideal functionalities that our constructions assume for secure
MPC over $\Zp$ and over a generic group $\G$. We benchmark two concrete
frameworks that realise these functionalities: ATLAS, in the honest-majority
setting, and MASCOT, in the dishonest-majority setting.

\subsection*{ATLAS}\label{sec:ATLAS}

ATLAS is a secret-sharing based MPC scheme designed for the honest-majority
setting~\cite{C:GLOPS21}. Because it operates with an honest majority of
participants, ATLAS achieves unconditional security against both
honest-but-curious and malicious adversaries.

ATLAS builds on an existing protocol, improving its round complexity by a
factor of two~\cite{C:DamNie07}. Since it operates in the honest-majority
setting, no public key operations are needed, and this setting additionally
lets the protocol achieve fairness and guaranteed output delivery against a
semi-honest adversary.

The starting point for multiplication is as follows. For shares $[a], [b]$,
parties locally multiply their shares, giving a degree-$2t$ sharing of the
result. $2t$ honest parties then reshare this as a degree-$t$ sharing, for use
in the next round of computation. In this way, the degree of the shares does
not grow exponentially with each layer of computation, and remains below the
total number of parties $n$ throughout.

In the following, subscripts on shares denote the degree of the sharing. To
evaluate a multiplication gate, parties prepare a pair of random sharings
$([r]_t, [r]_{2t})$ of a random shared value $r$, where the first sharing has
degree $t$ and the second degree $2t$. To compute $[c]_t = [a \cdot b]_t$
from $[a]_t$ and $[b]_t$, a pair $([r]_t, [r]_{2t})$ is consumed as follows. A
party ${P_{\mathrm{king}}}$ is chosen, and all parties locally compute $[e]_{2t} = [a]_t
\cdot [b]_t + [r]_{2t}$ and send these values to ${P_{\mathrm{king}}}$, who reconstructs
$e$ and distributes it back to all parties. Every party then locally computes
$[c]_t = e - [r]_t$.

ATLAS relies on the observation that when ${P_{\mathrm{king}}}$ is honest, and the sharing
of $e$ that ${P_{\mathrm{king}}}$ distributes is degree $t$ rather than degree $2t$, the
corrupted parties receive $t$ shares of a random degree-$t$ sharing $[e]_t$
from ${P_{\mathrm{king}}}$, which are uniformly random and independent of the secret. So
for an honest ${P_{\mathrm{king}}}$, the protocol for computing and consuming $([r]_t,
[r]_{2t})$ can be optimised, and the responsibilities of ${P_{\mathrm{king}}}$ are
themselves distributed across all parties. ATLAS reduces the number of rounds by half by combining ideas from Beaver
triples with existing techniques, letting parties use one round of
communication to cover two rounds' worth of multiplication
gates~\cite{C:DamNie07}. For the
second of these two layers, a Beaver triple $[x], [y], [x \cdot y]$ is
consumed rather than a pair $([r]_t, [r]_{2t})$; the triple is produced in the
pre-processing stage, and consumed in the online phase as follows. To
multiply $[a]$ and $[b]$, all parties combine $[a + x]$ to learn $a + x$, and
$[b + y]$ to learn $b + y$, and then compute $[a \cdot b] = (a + x)[a] - (b +
y)[b] + [x \cdot y]$. This lets two secret-shared values be multiplied in the
online phase, using random shared values prepared in the pre-processing
phase.

In terms of the functionalities of Section~\ref{sec:MPC}, since ATLAS is built
from a linear secret sharing scheme, addition can be performed by parties
simply adding their local shares. Multiplication is implemented as described
above, each instance consuming a pair $([r]_t, [r]_{2t})$, or a triple $[x],
[y], [x \cdot y]$, prepared during pre-processing. To combine shares, all
parties broadcast them to one another.

\subsection*{MASCOT}\label{sec:MASCOT}

MASCOT performs authenticated secure multiplication of shared field elements
in the maliciously secure dishonest-majority setting~\cite{CCS:KelOrsSch16}. Secret-sharing based MPC relies on generating multiplication
triples $[a], [b], [c] = [a \cdot b]$; the original SPDZ protocol, for the
dishonest-majority, malicious setting, generated these using somewhat
homomorphic encryption and zero-knowledge proofs.

MASCOT's contribution is to generate multiplication triples using oblivious
transfer (OT) instead, extended to the arithmetic setting. OT, at a high
level, involves a sender who inputs $s_0, s_1$ and a receiver who inputs a bit
$b$; the sender receives no output, and the receiver receives $s_b$, learning
nothing about $s_{1-b}$. MASCOT uses OT together with MACs, standard sacrifice
techniques, and privacy amplification, so that the protocol either correctly
outputs the product of the shared elements, or otherwise aborts once
malicious behaviour is detected.

Multiplication of secret-shared elements proceeds as follows. To generate
multiplication triple $[a], [b], [a \cdot b]$, one of the two elements to be
multiplied, say $[b]$, is split into its component bits. For elements of
bit-length $k$, $k$ OTs are then run in parallel: the sender inputs a random
value $r$ and its own value to be multiplied, $[a]$; the receiver inputs the
bit decomposition of its input, $(b_1, \ldots, b_k)$, and receives either
$r_i$ or $r_i + a$ depending on the corresponding input bit. The values
received satisfy $q_i = r_i + b_i \cdot a$, which, mapped back to a field
element, gives an additive sharing $q + r = a \cdot b$.

Extending this two-party construction to $n$ parties is done by having each
pair of parties run the protocol between themselves, and then summing the
resulting partial shares. In the malicious setting, share values are
accompanied by MACs, under a global MAC key that is itself secret-shared
across all parties; the MAC values are checked with every multiplication, and
a malicious party supplying inconsistent inputs causes the MAC check, and so
the protocol, to abort.

In terms of the functionalities of Section~\ref{sec:MPC}, since MASCOT is
built over a linear secret sharing scheme, addition of secret-shared values
happens locally, by parties adding the shares of the values along with their
corresponding MACs. Multiplication is as described above. Inversion of a
secret-shared value can also be performed locally, by each party inverting
its own share. Combining shared values is performed by broadcasting shares,
which are reconstructed through addition, with the MACs checked to ensure
the reconstruction is correct.

\paragraph{Implementation.} Our implementation builds on MP-SPDZ, a framework
for implementing MPC with pre-processing~\cite{CCS:Keller20}. MP-SPDZ is built for the paradigm of sharing elements of
$\Zp$, with all operations defined between elements of $\Zp$. For the
protocols of this chapter, we additionally need operations between elements
of $\group{1}$ and $\group{2}$. Extending MP-SPDZ with group operations,
concretely elliptic curve group operations, has been considered and
implemented in prior work~\cite{ESORICS:DOKSS20,NDSS:GPVRNC23}.

We need the groups used to be pairing-friendly. Although we do not distribute
any pairing operations, since these occur only in verification, we still need
all operations to take place over pairing-friendly curves. The RELIC toolkit
has been integrated into MP-SPDZ for operations over pairing-friendly curves
and shared computation over elements of the resulting elliptic curve
groups~\cite{NDSS:GPVRNC23,relictoolkit}.

MP-SPDZ operates with a pre-processing (offline) phase and an online phase.
In our setting, $\algo{OffSign}$ corresponds to the pre-processing phase, and
$\algo{OnSign}$ and $\algo{Combine}$ run in the online phase.

We benchmark both the honest-majority and dishonest-majority settings, using
ATLAS and MASCOT respectively to realise multiplication of shared secrets
~\cite{CCS:KelOrsSch16,C:GLOPS21}. ATLAS is designed to work in both the
semi-honest and malicious settings, but the implementation available in
MP-SPDZ captures only the semi-honest case~\cite{CCS:Keller20}. ATLAS and
MASCOT both implement the ideal functionality of Section~\ref{sec:formalization}:
MASCOT via the arithmetic black box functionality, which defines
\texttt{init}, \texttt{input}, \texttt{add}, \texttt{multiply}, and
\texttt{output}, and ATLAS via its own ideal functionality, which defines
similar operations. In both cases, $\idfu{\algo{Inv}}$ is built from $\idfu{\algo{Mul}}$, with
$\idfu{\algo{RandShare}}$ taking the role of \texttt{input} and $\idfu{\algo{Combine}}$ taking
the role of \texttt{output}.

We run benchmarks on a 16-core machine running RHEL 9.0-64 with 32GB of RAM,
inside Docker containers based on Ubuntu 22.

\begin{table*}[!ht]
  \centering
  \begin{tabular}{ c || c | c | c || c | c | c }
   \toprule
    & Assumption & Message space & $|\sk|$ & $|\pk|$ & $|\signature|$ & Verification cost \\
   \midrule
    ~\cite{AC:CKPSS23} & PS + ROM & $(\gen^{\vec{m}}, \ggen^{\vec{m}}) \in \group{1}^l \times \group{2}^l$ &
    $ (l+1) | \Zp | $ & $ l | \group{1} | + (l+1) | \group{2} | $ & $ 2 | \group{1} | $ & $3l+2$ \\
    ~\cite{PKC:MMSST24} & SXDH & $ \vec{M} \in \group{1}^l$ &
    $ 2(l+1) | \Zp | $ & $ (l+1) | \group{2} | $ & $ 6 | \group{1} | + | \group{2} | $ & $l+11$ \\
    Section~\ref{sec:tsps-abe} & q-type & $ \vec{M} \in \group{1}^l$ &
    $ (l+1) | \Zp | $ & $ | \group{1} | + l| \group{2} | $ & $ | \group{1} | + 2 | \group{2} | $ & $l+4$ \\
    Section~\ref{sec:tsps-kiltz} & SXDH & $ \vec{M} \in \group{1}^l$ &
    $ l |\Zp | + 6 | \group{1} | $ & $ (l+6)| \group{2} | $ & $ 6 | \group{1} | + | \group{2} |$ & $2l+10$ \\
   \bottomrule
  \end{tabular}
  \caption{Comparison of threshold structure-preserving signature schemes.
  Verification cost lists the number of pairing evaluations needed during
  verification.}
  \label{table:comp}
\end{table*}

\begin{table*}[!ht]
  \centering
  \begin{tabular}{ c || c | c | c }
   \toprule
    & $\algo{OffSign}$ & $\algo{OnSign}$ & $\algo{Combine}$ \\
   \midrule
    ~\cite{AC:CKPSS23} & 0 & 0 & $ 2 | \group{1} |$ \\
    ~\cite{PKC:MMSST24} & 0 & 0 & $ 6 | \group{1} | + | \group{2} |$ \\
    Section~\ref{sec:tsps-abe} & $ n(l+4)|\Zp|$ & 0 & $| \group{1} | + 2 | \group{2} |$ \\
    Section~\ref{sec:tsps-kiltz} & $ n(12|\Zp| + 6|\group{1}|) $ & 0 & $6 | \group{1} | + | \group{2} | $ \\
   \bottomrule
  \end{tabular}
  \caption{Per-party communication complexity for the signing and combining
  phases of the threshold signing schemes.}
  \label{table:comms}
\end{table*}

\begin{table*}[!ht]
  \centering
  \begin{tabular}{ c || c || c | c || c | c }
   \toprule
    & Centralised & \multicolumn{2}{c||}{ATLAS~\cite{C:GLOPS21}} & \multicolumn{2}{c}{MASCOT~\cite{CCS:KelOrsSch16}} \\
    & & $\algo{OffSign}$ & $\algo{OnSign}$ & $\algo{OffSign}$ & $\algo{OnSign}$ \\
   \midrule
    Section~\ref{sec:tsps-abe} & 10ms & 46ms & 56ms & 92ms & 84ms \\
    Section~\ref{sec:tsps-kiltz} & 40ms & 88ms & 61ms & 352ms & 94ms \\
   \bottomrule
  \end{tabular}
  \caption{Timings of the sub-protocols of our threshold structure-preserving
  signatures, averaged over 1000 runs, for $n = 5$.}
  \label{table:benchmarks}
\end{table*}

Table~\ref{table:comp} compares our two schemes against two prior threshold
structure-preserving signature schemes in terms of assumptions, message
space, key and signature size, and verification
cost~\cite{AC:CKPSS23,PKC:MMSST24}, while Table~\ref{table:comms} compares
per-party communication complexity across the offline, online, and combining
phases. Table~\ref{table:benchmarks} reports concrete timings for both of our
schemes, using ATLAS and MASCOT to realise the underlying MPC
functionalities, and against a centralised (non-threshold) baseline, for $n =
5$ parties.

\paragraph{Comparing across total number of parties.} We next examine how
execution time varies with the total number of parties, fixing the threshold
of participating parties and the message length at two group elements. In
each case, ATLAS is around twice as fast as MASCOT, which is expected, since
MASCOT relies on public key cryptography to handle dishonest majorities.
Figures~\ref{eval-keygen} and~\ref{eval-offline} further show that ATLAS
scales better than MASCOT as the number of parties increases. Since our
benchmarks build on MP-SPDZ, they reflect the standard MPC offline and online
phases, where the output is constructed within the online phase; this means
the online and combining phases are benchmarked together, so their timings
grow with the number of parties. In practice, when using threshold
signatures, parties may instead perform only the constant-time $\algo{OnSign}$
locally, and submit their signature shares to a single party to perform
$\algo{Combine}$, rather than all parties communicating shares and each
constructing a copy of the signature locally. The signature generation phase
shown in Figure~\ref{eval-online} combines $\algo{OnSign}$ and $\algo{Combine}$.
